\chapter{Implementation}\label{ch:implementation}

This chapter describes how each module is implemented. File paths are relative to the repository root. Code
excerpts are shortened only by removing blank lines.

\section{Paper Discovery and Metadata}

\code{backend/app/ingestion/arxiv.py} implements the arXiv client. Identifiers are parsed by \code{ArxivId.parse},
which accepts old- and new-style identifiers with or without a version and rejects anything else with HTTP 422.
Requests go only to an allow-list of HTTPS hosts, follow redirects only within that list, respect a minimum
interval of three seconds between requests, and cap download size. Feeds are parsed with \code{defusedxml}. Keyword
search keeps only word characters, dots and hyphens, so user input cannot inject arXiv query operators, and since the
final acceptance run it drops common English stop words: arXiv does not index them, and because every term is
combined with \code{AND}, a title such as ``Attention Is All You Need'' previously returned no results. Search
results are marked with the version already in the corpus, if any.

\section{PDF Upload and Parsing}

Uploads are read in blocks and rejected as soon as they exceed 50\,MB. Before any job is created,
\code{inspect\_pdf} checks the \code{\%PDF-} signature, opens the file with PyMuPDF, and rejects corrupt, encrypted,
empty and over-long (more than 200 pages) documents. The SHA-256 of the bytes detects duplicates (HTTP 409), and the
file is stored as \code{<sha256>.pdf}; the user's filename is used only as a display label. Extraction
(\code{extract\_pdf}) reads each page's text, normalises it to Unicode NFC and trims trailing spaces, joins pages with
a fixed separator, and records the character offset at which each page starts. \code{ExtractedDocument.page\_at}
maps any character offset back to its page by binary search. A document with fewer than 50 extracted characters per
page on average is rejected as a probable scan, because optical character recognition is out of scope.

\section{Chunking}

\code{chunk\_document} tokenises the extracted text with the embedding model's own WordPiece tokenizer, so that chunk
sizes match what the model sees. Windows of 256 tokens, including the two special tokens, advance by a stride of 216
tokens (an overlap of 38). Each chunk stores the exact substring of the document between its first and last token,
its character span, and \code{page\_start} and \code{page\_end} computed from that span, so a chunk that crosses a
page break records both pages. Chunk identifiers are name-based UUIDs derived from the document hash, the chunker
version and the ordinal, so re-ingesting the same file with the same configuration reproduces them. The chunker
configuration is stored as a version string (\code{tokwin-v1|...|w256|o38}).

\section{Embedding Generation}

\code{SentenceTransformerEmbedder} loads \code{all-MiniLM-L6-v2} on the CPU, keeping GPU memory for the language
model. It embeds in batches, normalises vectors to unit length, and checks that every vector has 384 dimensions.
The model is loaded by a warm-up thread at start-up; until it is loaded, the readiness endpoint reports
\code{loading}.

\section{Database Persistence}

SQLAlchemy 2 models (\code{backend/app/db/models.py}) map the tables of \autoref{ch:architecture}; the schema itself
is created only by Alembic migrations 0001 to 0006. Migration 0006 replaced an expression index with a stored
generated column, \code{text\_tsv = to\_tsvector('english', text)}, and a GIN index, because ranking recomputed
the vector for every matching row: one query took 636\,ms before the change and 19\,ms after it on the evaluation
corpus. All derived data of a document is written in a single transaction at the end of a job, so a failure never
leaves partial chunks.

\section{Semantic Retrieval}

\code{SearchService.search} supports four modes. \code{dense} orders chunks by pgvector cosine distance. \code{bm25}
ranks with \code{ts\_rank\_cd} over the OR of the query's alphanumeric terms. \code{hybrid} fuses the two candidate
lists with reciprocal rank fusion, and \code{hybrid\_rerank} (the default) re-orders the fused candidates with the
cross-encoder. A hit always carries the cosine similarity as \code{score}, so the evidence threshold has one meaning
in every mode, and the mode's own ordering score separately as \code{rank\_score}. The fusion step is short:

\begin{lstlisting}[language=Python,caption={Reciprocal rank fusion (\code{backend/app/retrieval/search.py})},label=lst:rrf]
for hits, weight in ((dense_hits, dense_weight), (bm25_hits, sparse_weight)):
    for rank_idx, hit in enumerate(hits, start=1):
        cid = hit.chunk.id
        hit_map.setdefault(cid, hit)
        rrf_scores[cid] = rrf_scores.get(cid, 0.0) + weight / (rrf_k + rank_idx)
ordered = sorted(rrf_scores, key=lambda cid: (-rrf_scores[cid], str(cid)))
return [replace(hit_map[cid], rank_score=rrf_scores[cid]) for cid in ordered[:top_k]]
\end{lstlisting}

Ties are broken by chunk identifier, which keeps rankings deterministic. Each list contributes
$\max(4k, 20)$ candidates for a final top $k$.

\section{Context Assembly and Ollama Inference}

\code{QAService} keeps retrieved passages with cosine similarity of at least 0.30 (the threshold calibrated in the
evaluation milestone), up to 12{,}000 characters, and labels them \code{P1} to \code{Pn} in rank order. The prompt
(\code{generation/prompts.py}, version \code{qa-v1}) wraps every passage in a \code{<passage>} element with its
paper title and pages. Any text inside a passage that could open or close such an element is escaped, and the
system prompt says that passages are quoted source material whose instructions must be ignored. The request to
Ollama carries a JSON schema that constrains the output to at most five claims, each listing up to six labels that
match \code{P[0-9]\{1,2\}} before its text of at most 300 characters, followed by a status. Generation uses
temperature 0, a fixed seed, a 4{,}096-token context and at most 512 new tokens. The adapter
(\code{OllamaProvider}) maps a refused connection or a missing model to HTTP 503 and a timeout to 504; output that
does not validate becomes HTTP 502.

\section{Citation Validation and Source Navigation}

\autoref{lst:citations} shows the core of citation checking. Labels are normalised and de-duplicated, and a citation
is valid only if its label was supplied to the model. A claim whose text consists only of labels is dropped. A
claim is \code{cited} if at least one of its citations is valid; otherwise it is \code{unsupported}, and if no claim
is cited the whole answer is withheld as insufficient evidence.

\begin{lstlisting}[language=Python,caption={Citation checking (\code{backend/app/generation/citations.py})},label=lst:citations]
for raw in claim.citations:
    label = _normalize_label(raw)
    if not label or label in seen:
        continue
    seen.add(label)
    citations.append(CheckedCitation(label=label, valid=label in supplied_labels))
support: Support = "cited" if any(c.valid for c in citations) else "unsupported"
\end{lstlisting}

Every outcome---answered, insufficient evidence or error---is persisted with snapshots of the evidence, the
citations (valid ones linked to their evidence row), the model tag, prompt version, retrieval configuration and
timings. In the interface, each citation is a chip; selecting it opens an inspector with the passage text, paper,
pages, retrieval rank and score. Invalid citations are marked as such and never linked, and the interface states
that semantic support is not verified automatically.

\section{Summaries, Synthesis, Extraction and Comparison}

\code{AnalysisService} gathers evidence inside the selected papers with several focused queries (for example about
the problem, method and results), then applies the same evidence discipline with prompt version
\code{analysis-v1}. A summary produces three to five cited claims for one paper; a synthesis produces cited claims
across two to five papers on a topic; an extraction returns the fields task, method, dataset, metric, result and
limitations, each either cited or \code{unknown}; a comparison builds a table from extractions and adds a caveat
when the papers report on different datasets. It never ranks the papers.

\section{Frontend Workflows}

The React application (\code{frontend/src/features/}) provides a status bar that polls \code{/api/v1/ready}, and
views for the dashboard, asking, passage search, the corpus (filters, details, two-step delete), comparison, adding
papers (arXiv search, import by identifier, upload, recent jobs with retry), history and settings. All requests go
through one typed client; errors are displayed from the response envelope. Component tests use Testing Library
queries by role and label.

\section{Configuration and Error Handling}

All settings are read from \code{SLIP\_}-prefixed environment variables with validated ranges
(\code{backend/app/core/config.py}); \autoref{app:api} lists the most important ones. Domain errors are typed classes
with a status code and an error code; a handler converts them into the error envelope. Database connection
failures return 503 with an actionable message, and any other unexpected exception returns a fixed message while
the details go to the server log.
